\documentclass[fontsize=14pt]{mybib}

\newcommand{\Name}{Philipp}
\setheaderlogo{mnr-green}

% ensure scrlayer-scrpage has sufficient footheight
\renewcommand{\DFA}{12mm}
\renewcommand{\DZA}{1.2}
\setlength{\footheight}{30.4pt}
\author{OTS}

\begin{document}
\begin{bibelbox}{SCHL}{Kol}{3:16-17}
    Lasst das Wort des Christus reichlich in euch wohnen in aller Weisheit; lehrt und ermahnt einander und singt mit Psalmen und Lobgesängen und geistlichen Liedern dem Herrn lieblich in eurem Herzen. Und was immer ihr tut in Wort und Werk, das tut alles im Namen des Herrn Jesus und dankt Gott, dem Vater, durch ihn.
\end{bibelbox}
\section{Begrüssung}
Ich möchte euch alle recht herzlich zum Gottesdienst begrüssen. Schön, dass ihr gekommen seid, um unserem \herr N zu danken, ihn zu loben und zu preisen. 
Hallo \Name{} schön, dass du wieder bei uns bist, um uns das Wort des Herrn zu bringen.

\begin{itemize}
    \item \beten[b, p]{Beten zur Begrüssung}
    \item und anschliessend singen wir zusammen das erste Lied
    \item \lied[b]{Gib Gott die Ehre}
\end{itemize}

\section{Informationen}
\begin{itemize}
    \item \info[b, p]{Bibel und Gebetsabend:} Do., 05. Nov. 2026 Bibel und Gebetsabend Nathanel Winkler Mk 11 1 - 11
    \item \info[b, p]{Nächster Gottesdienst:} So., 8. Nov. 2026 Gottesdienst  mit Nathanel Winkler
    \item \info[b, p]{Bibelleseabend:} Do., 12. Nov. 2026 19:00 Uhr Johannes Evangelium
    \item \info[b, p]{Adventsmarkt:} Do. 26., Nov. bis So. 29., Nov.2026 
    \item \betonung[b, p]{Kollekte:} Die Kollekte wird hinten in der grünen Box gesammelt und vollumfänglich für den Bau dieser Gemeinde eingesetzt.
\end{itemize}
\section{Einleitung}
Heute feiern wir unser Jahrestag hier in diesen Räumen. 
Armin und Marie-Therese haben beim Nachhause gehen vom Gottesdienst im Studio Fux, gesehen, dass hier unsere Vorgänger ausziehen.

Von da an ging es ziemlich schnell. 13.09.2024 kam Fredy von Spiez nach Naters, Fredy war zu der Zeit auf seiner Schweizer Tour. Fredy hat mit Herr D'Alpaos die Bedingungen ausgehandelt.  Herr D'Alpaos war zuerst etwas skeptisch, weil er keine Sekte in diesen Räumen haben wollte. Aber als Fredy ihm dann erklärt hat wer und was der MNR ist, hat er dann zugesagt. Nach dem Unterschreiben des Vertrages ging es an die Arbeit.

Danach ging es ans Teppich bestellen und Inventar besorgen. Ende Oktober kam dann noch ein Klavier. Alles viel Arbeit. Da möchte ich Oswald und Ariane danken. Sie haben viele Stunden in diese Räume investiert. Wie bei jedem Bau kam der Teppich sehr spät und am Ende war dann alles ein bisschen hektisch. Aber es hat geklappt.

Der erste Gottesdienst in diesen Räumen fand am ersten Sonntag im 3. November 2024 statt. Es gab einen sehr schönen Eröffnungsgottesdienst mit Fredy, Jonathan Malgo und Hovert Rose.

Diese Eröffnungsfeier, hat mir so richtig vor Augen geführt, dass dieses Projekt von Gott getragen wird. Das Lokal, der Umzug all das zeigt die Hand unseres geliebten Herrn. Von diesem Sonntag an wusste ich, das ist richtig. Darum wollen wir am ersten Sonntag im November diesen Tag feiern nicht, um uns auf die Schultern zu klopfen, sondern um Gott zu danken, was er hier aufgebaut hat. So wollen wir weiterhin für sein Reich hier in Naters Salz und Licht  sein.

Dann ab Mai 2025 zogen die ersten Wolken auf. Leider konnten die Differenzen zwischen mir, dem MNR und der Fam. Zuber nicht behoben werden. Die Familie Zuber hat sich im September 2025 zurückgezogen.

Mit der grossen Unterstützung der Familien Toth, Perren und Giz haben wir dann das Lokal neu eingerichtet und so konnten die Gottesdienste ohne Unterbrechung weiter gehen.

Highlights waren auch die zwei Gemeindefeste jeweils im August 25 und 26. Es war eine sehr schöne Gemeinschaft. Das aber ist Geschichte und eine von Gott geführte Geschichte. Ab Heute starten wir in das Abenteuer 2027. 
Mein Plan, 
\begin{enumerate}
    \item die Gottesdienste und Gebetsabende gehen weiter, wie bis jetzt. Jeweils 14:45Uhr und Do 20:00Uhr
    \item erfolgreicher Adventsmarkt
    \item Bibelverteilaktion
    \item Bibellesen mit plus Personen
    \item Gemeinde Fest im August
    \item ...
\end{enumerate}
Das alles funktioniert nur, wenn uns Gott dabei unterstützt. Solange wir uns an sein Wort halten, können wir nichts falsch machen, egal wie das Resultat ausfällt. So, dass wir am Ende wie Paulus sagen können: \enquote{Ich habe den guten Kampf gekämpft, den Lauf vollendet, den Lauf vollendet den Glauben bewahrt!}

Im Anschluss werde ich nochmals die Folie laufen lassen, die schon gesehen habt.
\begin{itemize}
    \item \lied[b]{Du bist unsere Zuversicht}.
\end{itemize}
\section{Predigt}
    \Name
\section{Abendmahl}
    Mit \Name
\section{Abschluss}
\begin{itemize}
    \item \lied[b]{Ich weiss nicht, warum Gottes Gnad}.
    \item \beten[b]{Zum Abschluss}
\end{itemize}
\begin{bibelbox}{SCHL}{IIKor}{13:13}
    Die Gnade des Herrn Jesus Christus und die Liebe Gottes und die Gemeinschaft des Heiligen Geistes sei mit euch allen! Amen
\end{bibelbox}
\end{document}